\section{Analysis}
\label{sec:analysis}

We state three results. The first bounds the resident set in sessions and shows why a
concurrency threshold cannot be the right control variable. The second prices a revival
and locates the real scarcity. The third bounds how much a wrong growth reserve can move
the ledger, which is the formal reason the design can afford a crude estimate.

\subsection{Prop.~1: the feasible set is small, and a threshold cannot track it}

Let the live sessions' footprints be $\Phi_{(1)}\le\cdots\le\Phi_{(m)}$ in ascending order
and let $\rho C_d$ be the committed device budget. A resident set is feasible exactly when
its members are the smallest footprints that fit:
\begin{equation}
|\mathcal{F}(t)|=\max\Bigl\{\,k:\sum_{i\le k}\Phi_{(i)}\le\rho C_d\,\Bigr\},
\label{eq:prop1}
\end{equation}
which is a statement about the multiset of footprints, not a mean-field quantity. On our
trace the count has median $1$ and mean $1.77$ over the 81 rounds with live sessions; at
63\% of rounds at most one session fits, and three or more fit at only 16\%. The scalar
estimate $\rho C_d/\bar\Phi=1.40$, with $\bar\Phi=69.0$K the mean session footprint, lands
in the same range but hides the skew: early rounds fit up to eight small sessions and late
rounds fit one. In the mean-field sense---working set $m$ times the pool---the resident
count is $|\mathcal{S}|/m=19/12.9=1.47$, which is the number an operator would size for.

Three consequences contradict a habit of the literature.
(i) A concurrency threshold cannot track this. A window of $k=2$ is infeasible at 63\% of
the rounds and $k=3$ at 84\%, while at the start of the trace a window of eight would have
been feasible; any fixed $k$ is either over-committed for most of the run or leaves most of
the card idle at the start.
(ii) Queueing reasoning about arrival rates is the wrong instrument: there is no stable
operating point with many sessions in flight, so latency is set not by utilisation but by
\emph{which} session holds the card.
(iii) The interesting quantity is not how many sessions are admitted but how much
recomputable work is kept alive per byte held---which is what \S\ref{sec:m4} ranks.

\subsection{Prop.~2: revivals are cheap; slots are not}

With a linear prefill model $P(c)=c/R$ and per-token KV size $\kappa$, substituting
\eqref{eq:trev} and \eqref{eq:trec} gives
\begin{equation}
\frac{\Trev(s)}{\Trec(s)}
=\frac{\kappa R}{B}+\frac{\tau R}{c}.
\label{eq:ratio}
\end{equation}
On our deployment $\kappa=56$\,KiB/token (the model's KV shape) and $R=465$\,tok/s (the
median over 131 cold rounds), with $B=\pend{GB/s}$ measured for host-to-device copies, so
the first term is $\pend{ratio}$ and \emph{restoring a session's state costs about two orders of magnitude
less than recomputing it}, for any context beyond a few hundred tokens.

This changes what the problem is. If revivals were expensive, the right design would be a
careful oracle about when to keep state; the measurement says the transfer is not the
bottleneck---the \emph{slot} is. The host tier holds $C_h=203$K tokens, and
the working set is $12.9\times$ the device pool and $6.5\times$ the host tier,
so at most 2.9 sessions can be parked at once and all the others must lose
their state. The decision that matters is a replacement decision between sessions.

Its marginal condition follows from (2): holding $s$ rather than $t$ is worth it when the
prefill time kept alive per byte is larger,
\begin{equation}
\frac{p_s\,\Trec(s)}{\Phi_s}\;>\;\frac{p_t\,\Trec(t)}{\Phi_t},
\label{eq:rank}
\end{equation}
with $p_s$ the probability that the session returns before its state would be evicted
anyway. The controller instantiates \eqref{eq:rank} in an observable form: victims are the
youngest, and the youngest session is the one with the smallest invested context, hence the
smallest $\Trec$. An idle-time gate---suspend only after a session has been parked longer
than $t^\ast$, the time to serve one round of another session---is the natural refinement
of \eqref{eq:rank}, because a session that returns sooner than $t^\ast$ would pay a revival
for a slot the replacement could not have used. It is implemented but disabled in the runs
we report (\S\ref{sec:m4}), so the evaluation reflects youngest-first shedding alone.

Equation~\eqref{eq:rank} is a \emph{replacement} condition, not a per-session one: when the
host tier is full, suspending a session always ejects another, and the cost of a
suspension is the revival-or-recompute it causes elsewhere. This is why the design tracks
host occupancy instead of treating the host tier as free capacity.

\subsection{Prop.~3: a wrong growth reserve cannot move the ledger far}

Write the projection error of session $s$ as
$\Phi_s^{\text{next}}-\hat\Phi_s = g_s-\hat g_s$, with $g_s$ the realised growth. Then
\begin{equation}
\frac{|\Phi_s^{\text{next}}-\hat\Phi_s|}{\Phi_s}
\;\le\;\frac{g_s}{\Phi_s}+\frac{|\hat g_s-g_s|}{\Phi_s}
\;\le\;\varepsilon_s+\varepsilon'_s ,
\label{eq:err}
\end{equation}
where $\varepsilon_s=g_s/\Phi_s$ is the growth-to-footprint ratio and $\varepsilon'_s$
its estimation error, both normalised by the same footprint. On the trace the reserve
$\hat g_s$ has a median of 4.0K tokens against a 48.5K-token context, so
$\varepsilon_s\approx8.2\%$, and the margin of \S\ref{sec:m3} makes the projection exceed
the realised footprint in all 583 round transitions, by a median of 22.1\% of the context.
The projection therefore errs in the safe direction, and its error is dominated by the
margin, not by the estimate: replacing the per-session EMA with the global mean growth
moves the projection by less than the margin does.

The cost of that conservatism is capacity. Sweeping the two constants offline over the same
583 transitions gives coverage between 67.4\% ($h=1$, no margin) and 100\%, with slack
between 0.2\% and 22.1\%; the deployed corner ($h=3$, margin $1.15$) is the most
conservative of the nine, and the cheapest corner that still covers 99.5\% of transitions
($h=2$, margin $1.05$) would hold back 8.0\% instead of 22.1\%. This is the one place where
the design leaves performance on the table, and it is a tuning question rather than a
structural one---which is why we report the sweep instead of claiming the setting is
optimal.

We report the empirical version of \eqref{eq:err} in \S\ref{sec:eval}: the distribution of
the realised projection error over all admissions, and the fraction of admitted sessions
whose next round exceeded the reserve.
